\section{Geometry of the actual composite ordinary constraint}
\label{cmg-section}

Fix an odd primitive ratio $\ell\ge3$, a positive integer $s$, and
$R=\ell^s$. Current spacing is one. An earlier classical refinement
$M=\ell^k$ does not enter any map in this section. Constants are uniform
in $M$ and compatible periodic volume once $\ell,s$ are fixed; no
uniformity in $R\to\infty$ is asserted. Compatible source periods
$2\ell^n$, $n\ge s$, suffice. Local paths are evaluated in their specified
lifts, so the two opposite interfaces of a two-block periodic direction
remain distinct. Translate the common block center to zero and identify
the level-$t$ lattice with $\ell^t\mathbb Z^d$.

The constraint constructed here is the actual composition of $s$
ordinary primitive averages. It is not identified with a direct-$R$
ordinary average or with $s$ separately truncated quantum integrals.
The direct-$R$ contour variable used below is a separate gauge choice.

\subsection{The source map and its exact covariance}
\label{cmg-source-section}

For a centered primitive block and $x\in B_\ell(y)$, let
$\mathcal G_\ell(y,x)$ contain the shortest coordinate-order paths
from $y$ to $x$. Every such path is internal to this block. Use the
selected source group mean $\mathfrak M$ to define
\begin{equation}
 P_\ell(V;y,x)=
 \mathfrak M\bigl(\{V(\Gamma):\Gamma\in\mathcal G_\ell(y,x)\}\bigr).
 \label{cmg-contour}
\end{equation}
For $c=(y,y+\ell e_\mu)$, put $x'=x+\ell e_\mu$ and let
$R_c(V)$ be the straight center transporter. The ordinary map is
\begin{align}
 U_x(V)&=P_\ell(V;y,x)V([x,x'])P_\ell(V;y+\ell e_\mu,x')^{-1},
 \notag\\
 \mathcal M_\ell(V)(c)
 &=\exp\left\{\ell^{-d}\sum_{x\in B_\ell(y)}
          \log[U_x(V)R_c(V)^{-1}]\right\}R_c(V).
 \label{cmg-ordinary}
\end{align}
All logarithms use one common small branch. These are the ordinary
average and contour definitions in \cite{balaban109cm}, pp.~252--254,
(0.3), (0.5)--(0.12).

The source axioms give local analyticity, the arithmetic differential
at equal identity inputs, and endpoint covariance
\begin{equation}
 \mathfrak M(\{gU_jh^{-1}\})=
 g\mathfrak M(\{U_j\})h^{-1}.
 \label{cmg-mean-covariance}
\end{equation}
For completeness, (0.10) of that paper supplies one admissible
realization as the unique small solution $P$ of
$\sum_j\log(U_jP^{-1})=0$. The derivative with respect to $\log P$ at
identity is minus the number of inputs times the identity, so the
analytic implicit-function theorem gives uniqueness, analyticity and
the arithmetic differential. Conjugating the defining logarithms
proves \eqref{cmg-mean-covariance} for this realization. The paper
explicitly permits any mean satisfying its axioms; the proof here
uses those properties for the selected mean and does not replace it
by a different one.

With $V^u(b)=u(b_-)V(b)u(b_+)^{-1}$, it follows that
\begin{align}
 P_\ell(V^u;y,x)&=u(y)P_\ell(V;y,x)u(x)^{-1},\notag\\
 \mathcal M_\ell(V^u)&=
       \mathcal M_\ell(V)^{\,u|_{\text{new centers}}}.
 \label{cmg-primitive-covariance}
\end{align}
Indeed every relative transporter inside the logarithm in
\eqref{cmg-ordinary} is conjugated at its source center. The identities
hold on the common analytic domains and extend to complex fields there.

\subsection{One primitive pivot has one nonlinear consumer}
\label{cmg-primitive-section}

Select the unit interface bond
\begin{equation}
 p_\ell(c)=
 [y+(\ell-1)e_\mu/2,\ y+(\ell+1)e_\mu/2].
 \label{cmg-primitive-pivot}
\end{equation}
No internal contour reads this bond. A parallel string containing it
must have direction $\mu$, the same transverse block center, and
the same pair of longitudinal endpoint blocks. Hence it belongs to
the output equation $c$. The reference center string has the same
property. This is a statement about the actual products and group
means before differentiation:
\begin{equation}
 \mathcal M_\ell(V)(c')\text{ is independent of }V(p_\ell(c))
 \quad(c'\ne c).
 \label{cmg-primitive-incidence}
\end{equation}
The local lifted strings distinguish opposite periodic interfaces.

Let $C_\ell$ be the angle differential at identity, and let
$t_y(B;x)$ be the arithmetic mean of the oriented sums of $B$ on
the root-to-$x$ contours. Linearizing \eqref{cmg-ordinary} gives
\begin{equation}
 C_\ell B(c)=\ell^{-d}\sum_{x\in B_\ell(y)}
 \left[t_y(B;x)+B([x,x+\ell e_\mu])
             -t_{y+\ell e_\mu}(B;x+\ell e_\mu)\right].
 \label{cmg-primitive-linear}
\end{equation}
The reference-string derivative cancels between the logarithms and
the final factor. The root terms do not read the pivot. Exactly $\ell$
starting sites have its transverse coordinates, and each corresponding
positive string contains it once. Thus its full column is
\begin{equation}
 (C_\ell)_{c',p_\ell(c)}
       =\ell^{1-d}\delta_{c,c'}I_{\mathfrak g}.
 \label{cmg-primitive-coefficient}
\end{equation}
This is an angle derivative, without division by the coarse length.

\subsection{Exact nonlinear incidence through the composition}
\label{cmg-composite-section}

Define
\begin{equation}
 \mathcal M_R^{\mathrm{comp}}=\mathcal M_\ell^s.
 \label{cmg-composite}
\end{equation}
For a final bond $c=(y,y+Re_\mu)$ choose
\begin{equation}
 p_R(c)=[y+(R-1)e_\mu/2,\ y+(R+1)e_\mu/2].
 \label{cmg-macro-pivot}
\end{equation}
Writing $h_t=\ell^t$, set
\begin{equation}
 b_t(c)=[y+(R-h_t)e_\mu/2,\ y+(R+h_t)e_\mu/2],
 \qquad 0\le t\le s.
 \label{cmg-pivot-chain}
\end{equation}
Its endpoints are level-$t$ sites because $R/h_t$ is odd.
Moreover $b_0=p_R(c)$, $b_s=c$, and $b_t$ is exactly the primitive
middle-interface pivot of $b_{t+1}$.

\begin{proposition}
\label{cmg-incidence-proposition}
Changing the input group coordinate at $p_R(c)$ changes at level $t$
at most $b_t(c)$. Consequently every final constraint reads its own
selected pivot and no other selected pivot. Each final output reads
only unit bonds whose endpoints lie in its two endpoint $R$-blocks,
using the prescribed local lift. If $C_R$ denotes the composite angle
differential at identity, then
\begin{equation}
 (C_R)_{c,p_R(c')}=R^{1-d}\delta_{c,c'}I_{\mathfrak g}.
 \label{cmg-composite-coefficient}
\end{equation}
\end{proposition}

\begin{proof}
The incidence assertion is induction along
\eqref{cmg-pivot-chain}, applying
\eqref{cmg-primitive-incidence} at spacing $h_t$.
It is functional independence, so is valid for nonlinear variations,
not only for the derivative. Distinct final bonds have distinct
intermediate selected bonds.

For locality, a level-$t$ output reads level-$(t-1)$ bonds in its two
endpoint primitive blocks, including their connecting strings.
The descendants centered at those intermediate sites tile the two
centered blocks of side $\ell^t$. Iterating gives the stated unit-bond
dependency set. On short compatible tori this is read in the local
cover, not inferred merely from periodic set containment.

Finally the chain rule along the single affected output multiplies
$s$ copies of $\ell^{1-d}I$. There is no intermediate row sum.
The product is $R^{1-d}I$, proving
\eqref{cmg-composite-coefficient}.
\end{proof}

Composition of \eqref{cmg-primitive-covariance} gives
\begin{equation}
 \mathcal M_R^{\mathrm{comp}}(V^u)=
 \mathcal M_R^{\mathrm{comp}}(V)^{\,u|_{R\text{-centers}}}.
 \label{cmg-composite-covariance}
\end{equation}
With $D\phi(b)=\phi(b_+)-\phi(b_-)$, differentiate using
$u=e^{-it\phi}$ to obtain
\begin{equation}
 C_RD\phi(c)=\phi(c_+)-\phi(c_-).
 \label{cmg-gradient}
\end{equation}
The right side is an endpoint difference, or $R$ times the physical
coarse gradient.

\subsection{Local analytic elimination and unchanged retained coordinates}
\label{cmg-chart-section}

For a small background $V=e^{ia}$ define
\begin{align}
 F_{R,V}(B;c)&=\frac1i\log\left[
 \mathcal M_R^{\mathrm{comp}}(e^{iB}V)(c)
 \mathcal M_R^{\mathrm{comp}}(V)(c)^{-1}\right],\notag\\
 A_{R,V}&=D_BF_{R,V}(0).
 \label{cmg-relative-map}
\end{align}
Finite composition and the two-block dependency give a positive
product domain in $a,B$, independent of volume and old $M$ at fixed
$R$. Put
\begin{equation}
 \pi_c(V)=R^{d-1}(A_{R,V})_{c,p_R(c)}.
 \label{cmg-pi}
\end{equation}
It is local analytic, $\pi_c(1)=I$, and on a common smaller domain
$\lVert\pi_c-I\rVert\le1/2$, $\lVert\pi_c^{-1}\rVert\le2$.
The right inverse inserts only at selected pivots:
\begin{equation}
 (h_{R,V}w)(p_R(c))=R^{d-1}\pi_c(V)^{-1}w(c),\qquad
 (h_{R,V}w)(b)=0\ (b\text{ nonpivot}),\qquad A_{R,V}h_{R,V}=I.
 \label{cmg-right-inverse}
\end{equation}
In particular its maximum row norm is at most $2R^{d-1}$.
If $Q_{R,V}=A_{R,V}/R$ is the physical coarse-potential derivative,
and its kernel is written as
$QB(c)=\sum_bR^{-d}Q(c,b)B(b)$, that kernel's pivot coefficient
is exactly $\pi_c(V)$. Thus the factors in
\eqref{cmg-right-inverse} retain the source physical normalization.

Let $E_f$ insert nonpivot values and zero at pivots, and abbreviate
$A=A_{R,V}$, $h=h_{R,V}$. Define
\begin{equation}
 K_{R,V}=E_f-hAE_f,\qquad AK_{R,V}=0.
 \label{cmg-insertion}
\end{equation}
Restriction of $K_{R,V}z$ to nonpivots is $z$. Write
$F_{R,V}(B)=AB+\mathcal C_V(B)$, with
$\mathcal C_V(0)=D\mathcal C_V(0)=0$, and solve
\begin{equation}
 \widetilde{\mathcal D}_V(B)
       =\mathcal C_V(B-h\widetilde{\mathcal D}_V(B)).
 \label{cmg-implicit-correction}
\end{equation}
For output $c$ the right side reads only its own unknown
$\widetilde{\mathcal D}_V(B;c)$: other unknowns insert at pivots that
the full nonlinear $c$ equation does not read. The implicit derivative
at zero is identity. A finite-dimensional analytic implicit theorem,
or contraction after imposing $\lVert D\mathcal C_V h\rVert<1/2$,
therefore gives a unique local analytic solution. There are finitely
many oriented neighborhood types, so one radius works in all volumes.
Iteration does not spread its spatial dependency.

The exact retained chart is
\begin{align}
 \Phi_{R,V}(z)
   &=K_{R,V}z-h\widetilde{\mathcal D}_V(K_{R,V}z),\notag\\
 F_{R,V}(\Phi_{R,V}(z))&=0,\qquad
 \Phi_{R,V}(z)(b)=z(b)\quad(b\text{ nonpivot}).
 \label{cmg-retained-chart}
\end{align}
It is the unique small solution with its prescribed retained values.
Direct differentiation gives
\begin{align}
 D\Phi_{R,V}(0)&=K_{R,V},\notag\\
 D^2\Phi_{R,V}(0)[v,w]
   &=-hD^2F_{R,V}(0)[K_{R,V}v,K_{R,V}w].
 \label{cmg-chart-derivatives}
\end{align}
The affine equation $F_{R,V}(B)=w$ is equally local:
\begin{equation}
 B=K_{R,V}z+hw-h\mathcal C_V(B).
 \label{cmg-affine-chart}
\end{equation}
It has a common positive small product domain in $a,w,z$ and still
retains $z$ exactly. Local finite dimension, Cauchy estimates and
bounded neighborhood incidence give finite weighted first-derivative
row and column bounds and finite contracted second-derivative bounds.
Their constants may depend on $R$, but not on $M$ or volume.

After coupling insertion the relative field is
$\Phi_{R,V}(gz)$, whose retained values equal $gz$.
It is not $g\Phi_{R,V}(z)$. Additional pivot, contour or reconstructed
field indicators remain at their exact nonlinear arguments.

\subsection{An independent direct-$R$ contour gauge}
\label{cmg-gauge-section}

Use the same selected source mean on root-to-site coordinate-order
paths inside one centered $R$-block:
\begin{equation}
 P_R(V;y,x)=
 \mathfrak M\bigl(\{V(\Gamma):\Gamma\in\mathcal G_R(y,x)\}\bigr).
 \label{cmg-direct-contour}
\end{equation}
This does not redefine $\mathcal M_R^{\mathrm{comp}}$. Endpoint
covariance gives
\begin{equation}
 P_R(V^u;y,x)=u(y)P_R(V;y,x)u(x)^{-1}.
 \label{cmg-direct-covariance}
\end{equation}
Every path is internal to its block; every selected pivot crosses a
block interface. Thus $P_R$ is exactly independent of all selected
pivot group variables at every allowed background.

Define
\begin{equation}
 T_{R,V}B(y,x)=
 D_B\left.\frac1i\log[
 P_R(e^{iB}V;y,x)P_R(V;y,x)^{-1}]\right|_{B=0}.
 \label{cmg-T-definition}
\end{equation}
The pivot support of $h_{R,V}$ then gives
\begin{equation}
 T_{R,V}h_{R,V}=0,\qquad
 T_{R,V}K_{R,V}=T_{R,V}E_f.
 \label{cmg-T-identities}
\end{equation}
More strongly,
\begin{equation}
 P_R(e^{i\Phi_{R,V}(z)}V;y,x)
       =P_R(e^{iE_fz}V;y,x).
 \label{cmg-gauge-retained}
\end{equation}
The contour penalty therefore reads exactly the unchanged retained
variables, although its coefficients still depend on $V$.

For normalized trace $\tau$ put
\begin{equation}
 \mathcal G_R(V)=\sum_y\sum_{\substack{x\in B_R(y)\\x\ne y}}
       \left[1-\tfrac12\tau(P_R(V;y,x)+P_R(V;y,x)^{-1})\right].
 \label{cmg-gauge-action}
\end{equation}
It is nonnegative on real unitary fields. At a background with every
$P_R=I$ its first derivative vanishes and its pulled-back Hessian,
in counting/trace coordinates, is
\begin{equation}
 D_z^2\mathcal G_R(e^{i\Phi_{R,V}(z)}V)\big|_{z=0}
       =K_{R,V}^{T}T_{R,V}^{T}T_{R,V}K_{R,V}.
 \label{cmg-gauge-hessian}
\end{equation}
The second derivative of the chart contributes nothing because it
is pivot supported and the gauge action is pivot independent. At
identity, \eqref{cmg-direct-covariance} also gives
\begin{equation}
 T_{R,1}D\phi(y,x)=\phi(x)-\phi(y).
 \label{cmg-T-gradient}
\end{equation}

\subsection{Clearing this gauge and the domain boundary}
\label{cmg-clearing-section}

Set $u_V(y)=I$ at all final centers and
$u_V(x)=P_R(V;y,x)$ elsewhere in its block. Then, on the common
domain,
\begin{equation}
 P_R(V^{u_V};y,x)=I,\qquad
 \mathcal M_R^{\mathrm{comp}}(V^{u_V})
       =\mathcal M_R^{\mathrm{comp}}(V).
 \label{cmg-clearing}
\end{equation}
This is a geometric gauge operation, not by itself a minimum or
stationarity argument.

There are at most $d!$ contour orderings and every path has length
at most $d(R-1)/2$. In a fixed matrix norm,
$\lVert\prod_b e^{ia_b}-I\rVert\le
\exp(\sum_b\lVert a_b\rVert)-1$. The selected local mean has a
fixed Lipschitz constant near equal identity inputs, with the number
of inputs bounded by $d!$. Consequently, when $R\lVert a\rVert_\infty$
is sufficiently small,
\begin{equation}
 \lVert\log u_V/i\rVert_\infty
       \le C R\lVert a\rVert_\infty.
 \label{cmg-gauge-size}
\end{equation}
The constant here depends on dimension, group and the selected mean,
not on $R$. For a pre-gauge representative
$a_p=O(w/R)$ this gives $u=O(w)$. It does not give a rooted link-angle
bound $O(w/R)$: endpoint gauge differences may have size $O(w)$.
Nor does \eqref{cmg-gauge-size} make the local elimination radius
uniform in $R$.

The composite constraint, its actual local chart, and the independent
contour gauge are thus established. Gaussian coercivity and locality,
the physical minimum comparisons, the exact compact gauge insertion
and all cutoff normalizations are separate modules. None of these
geometric identities supplies a repeated quantum RG induction.
